\chapter*{Abstract}
\vspace{-\baselineskip}
\noindent Computational notebooks support data science learning, yet practical assistance inside browser-based editors remains limited. Code assistants optimised for file-based repositories poorly serve Kaggle notebooks, where work proceeds through indexed cells in a dynamic web interface. This Final Year Project presents an agentic large language model assistant for Kaggle computational notebooks, implemented as a Chrome extension with a Python native host. The system captures notebook context via DOM-grounded scraping, maintains live and persistent JSON snapshots, and exposes tool-mediated cell insertion, editing, deletion, and execution. Three interaction modes---Ask, Code, and Agentic---enforce distinct contracts: read-only explanation, copy-paste code with placement guidance, and batched tool dispatch under a mandatory two-phase protocol. The host validates tool arguments, coerces session metadata, and applies deterministic ordering for index-sensitive multi-cell operations. Empirical evaluation on frozen Kaggle snapshots using GLM~4.7 reports Ask mode achieving 89\% workflow topic coverage with zero side effects and reliable cell-level explanation and error diagnosis; Code mode passing four scenarios with runnable Python cells and zero browser writes. The prototype demonstrates that effective in-browser notebook assistance is feasible without privileged platform APIs through constrained tool orchestration and grounded context extraction, contributing a practical blueprint for safe LLM-assisted workflows in web-native notebook environments.

\vspace{1.5ex}
\noindent\textbf{Keywords:} Computational notebooks, large language models, Kaggle notebooks, Chrome extension, native messaging, DOM grounding, agentic artificial intelligence, tool calling, data science, GLM-4.7
